\documentclass[11pt,a4paper]{article}
\usepackage[utf8]{inputenc}
\usepackage[T1]{fontenc}
\usepackage[brazil]{babel}
\usepackage[margin=2.2cm]{geometry}
\usepackage{booktabs}
\usepackage{tabularx}
\usepackage{listings}
\usepackage{xcolor}
\usepackage{hyperref}
\usepackage{parskip}

\lstset{basicstyle=\ttfamily\footnotesize,frame=single,breaklines=true,
  columns=fullflexible,backgroundcolor=\color{gray!8},rulecolor=\color{gray!40}}
\hypersetup{colorlinks=true,linkcolor=black,urlcolor=blue}

\title{\textbf{TP2 --- Orquestração K3s, Proxy Redundante (HA)\\e Persistência de Estado Não Destrutiva}\\[4pt]
\large DCC091A --- Tópicos em Redes e Processamento Distribuído I}
\author{Grupo 6}
\date{10 de novembro de 2026}

\begin{document}
\maketitle

\section{Objetivo}
Evoluir a infraestrutura do TP1 para um cluster K3s de dois nós protegido por dois proxies
Nginx em alta disponibilidade (Keepalived com VIP flutuante 192.168.56.254), com uma
aplicação web e um banco de dados persistente, e demonstrar que o ambiente pode ser
pausado (\texttt{vagrant suspend}/\texttt{halt}) e retomado sem reprovisionamento.

\section{Arquitetura}

\begin{table}[h]
\centering
\begin{tabularx}{\textwidth}{llllX}
\toprule
VM & Nome & IP & vCPU/RAM & Papel \\
\midrule
1 & bastion    & .10 & 1 / 1024 MB & Nó de controle Ansible, kubectl, build da imagem, Nginx + Keepalived MASTER (101) \\
2 & k3s-server & .20 & 2 / 2048 MB & K3s server: control plane, datastore, CoreDNS, local-path \\
3 & k3s-worker & .30 & 2 / 2048 MB & K3s agent \\
4 & auxiliary  & .40 & 1 / 1024 MB & Nginx + Keepalived BACKUP (100), backup em \texttt{./data\_backup} \\
\bottomrule
\end{tabularx}
\caption{Máquinas virtuais (Ubuntu 22.04, rede host-only 192.168.56.0/24, interface \texttt{enp0s8}).}
\end{table}

Fluxo de uma requisição: cliente $\rightarrow$ VIP 192.168.56.254:80 (proxy que detém o VIP)
$\rightarrow$ upstream \texttt{least\_conn} para 192.168.56.20:30080 e 192.168.56.30:30080
(NodePort) $\rightarrow$ kube-proxy $\rightarrow$ uma das duas réplicas da API $\rightarrow$
PostgreSQL (Service ClusterIP \texttt{postgres}).

\subsection{Aplicação}
API em Go (\texttt{app/main.go}) com as rotas \texttt{/info} (Pod e nó que atenderam, total
de registros), \texttt{/health} (200 se o banco responde, usada nas probes) e
\texttt{/records} (GET lista, POST grava). O Dockerfile é multi-stage: compila em
\texttt{golang:1.22-alpine} e copia apenas o binário estático para uma imagem
\texttt{FROM scratch}, executada como usuário não-root. A imagem final tem
\textbf{5,73 MB} (limite do roteiro: 100 MB).

A imagem é gerada com \texttt{podman} no bastion, exportada como \texttt{docker-archive}
e importada no containerd de cada nó (\texttt{k3s ctr images import}), com
\texttt{imagePullPolicy: Never}. Assim não é preciso um registry.

\subsection{Manifestos (\texttt{k3s/})}
\begin{itemize}
\item \texttt{01-configmap-secret.yaml}: ConfigMap \texttt{app-config} (ambiente, porta, host,
  nome e usuário do banco) e Secret \texttt{db-credentials} (senha).
\item \texttt{02-pvc.yaml}: PVC \texttt{postgres-data}, 1 Gi, \texttt{local-path}, RWO.
\item \texttt{03-deployment.yaml}: PostgreSQL (1 réplica, estratégia \texttt{Recreate}, PVC) e
  API (2 réplicas, \texttt{topologySpreadConstraints} por nó, liveness e readiness em
  \texttt{/health}, nome do Pod e do nó injetados pela Downward API).
\item \texttt{04-service.yaml}: \texttt{postgres} (ClusterIP) e \texttt{tp2-api} (NodePort 30080).
\end{itemize}

\subsection{Automação (Ansible a partir do bastion)}
O \texttt{Vagrantfile} define o bastion por último: ele recebe as chaves privadas das outras
VMs (\texttt{.vagrant/machines}) e executa o \texttt{site.yml} via \texttt{ansible\_local}.
\begin{itemize}
\item \textbf{common}: pacotes (curl, htop, tcpdump, dnsutils, jq), fuso
  America/Sao\_Paulo, \texttt{/etc/hosts} com todos os nós e o VIP, \texttt{ip\_forward} e UFW
  (entrada negada por padrão; regras por papel).
\item \textbf{proxy\_ha}: Nginx com upstream para os NodePorts, Keepalived com VRRP unicast,
  \texttt{check\_nginx.sh} e \texttt{ip\_nonlocal\_bind}.
\item \textbf{k3s\_server}: \texttt{/etc/rancher/k3s/config.yaml} (IP e flannel em
  \texttt{enp0s8}, Traefik e ServiceLB desativados), instalação na versão fixa
  v1.31.4+k3s1, extração do token e exportação do kubeconfig para o bastion.
\item \textbf{k3s\_agent}: instalação do agent e junção ao server com o token extraído.
\item \textbf{app\_deploy}: build da imagem (só quando o código muda), distribuição,
  \texttt{kubectl apply} e espera pelo rollout.
\item \textbf{backup\_sync} (extra): timer na VM 4 que grava \texttt{/records} em
  \texttt{/var/lib/app\_data} (synced\_folder \texttt{./data\_backup}) a cada minuto.
\end{itemize}

\section{Análise da proposta inicial}
A primeira versão do grupo acertava a topologia e os parâmetros, mas tinha problemas que
impediriam o funcionamento ou os testes, e que foram corrigidos:
\begin{itemize}
\item a interface \texttt{eth1} não existe na box \texttt{ubuntu/jammy64} (o nome é
  \texttt{enp0s8}), então o VIP e o flannel não subiriam;
\item o Ansible seria executado no host com a chave insegura global do Vagrant, mas o
  Vagrant gera uma chave por VM e o roteiro pede o Ansible no bastion;
\item a aplicação (\texttt{http-echo}) não tinha build multi-stage, banco de dados nem as
  rotas \texttt{/health} e \texttt{/info};
\item o PVC \texttt{local-path} RWO montado nas duas réplicas as prende ao mesmo nó;
\item o UFW era instalado sem regras, o kubeconfig não era exportado, e o
  \texttt{creates} do agent apontava para um arquivo inexistente (sem idempotência).
\end{itemize}

\section{Problemas encontrados durante a validação}
\begin{itemize}
\item \textbf{Ativação do UFW travava o playbook.} Na primeira ativação o conntrack passa a
  valer e a conexão SSH já aberta pelo Ansible nos nós K3s era descartada. A tarefa passou
  a rodar em \texttt{async}, seguida de \texttt{meta: reset\_connection} e
  \texttt{async\_status}.
\item \textbf{Failover em $\sim$4 s.} Com \texttt{weight -20} o master caía para prioridade
  81, mas o BACKUP descarta anúncios de prioridade menor e só assume após o
  \texttt{master\_down\_interval}. Sem \texttt{weight}, a falha do script coloca o master
  em FAULT, ele solta o VIP e anuncia prioridade 0: o tempo caiu para $\sim$1 s.
\item \textbf{Idempotência do Secret.} Com \texttt{stringData}, o \texttt{kubectl apply}
  sempre responde \texttt{configured}; o Secret passou a usar \texttt{data} (base64).
\item \textbf{Portas 80/443 abertas nos nós K3s} pelo Traefik e pelo ServiceLB embutidos,
  que não são usados; foram desativados.
\end{itemize}

\section{Resultados dos testes}
Todos os resultados abaixo foram obtidos em um ambiente recriado do zero
(\texttt{vagrant destroy -f \&\& vagrant up}, 11 min 46 s, sem falhas) em um host Windows 11
com VirtualBox 7.2.16 e Vagrant 2.4.9.

\begin{table}[h]
\centering
\begin{tabularx}{\textwidth}{lXl}
\toprule
Teste & Resultado & Status \\
\midrule
Idempotência & 2ª execução do \texttt{site.yml}: \texttt{changed=0, failed=0} nas 4 VMs & OK \\
A --- VIP e balanceamento & 10 consultas a \texttt{/info}: 6 na réplica do k3s-worker e 4 na do k3s-server & OK \\
B --- failover do proxy & Nginx parado às 13:44:49,40; bastion em FAULT às 49,88; VM 4 MASTER às 50,50 (\textbf{1,09 s}). 1 de 48 requisições perdida (a que estava em curso). VIP volta ao bastion quando o Nginx retorna & OK \\
C --- suspend/resume & API de volta pelo VIP em \textbf{51 s}; registro gravado antes preservado & OK \\
C --- halt/up & API de volta em 108 s; pods, VIP e registros restabelecidos sem reprovisionar & OK \\
D --- reconciliação & Pod apagado recriado, 2/2 réplicas prontas em \textbf{4 s} & OK \\
Backup (VM 4) & \texttt{records.json} e \texttt{backup.log} atualizados no host a cada minuto & OK \\
\bottomrule
\end{tabularx}
\caption{Resumo da validação.}
\end{table}

\subsection*{Trechos de saída}
Teste A:
\begin{lstlisting}
13:44:32 proxy=bastion {"node":"k3s-worker","pod":"tp2-api-7598cc9988-rpg9f",...}
13:44:33 proxy=bastion {"node":"k3s-server","pod":"tp2-api-7598cc9988-ht7g2",...}
13:44:34 proxy=bastion {"node":"k3s-worker","pod":"tp2-api-7598cc9988-rpg9f",...}
\end{lstlisting}

Teste B (logs do Keepalived):
\begin{lstlisting}
13:44:49.401 bastion   systemd: Stopped A high performance web server ...
13:44:49.884 bastion   Keepalived_vrrp: (VI_1) Entering FAULT STATE
13:44:50.496 auxiliary Keepalived_vrrp: (VI_1) Entering MASTER STATE
13:45:15.499 auxiliary Keepalived_vrrp: (VI_1) Entering BACKUP STATE
\end{lstlisting}

Teste C (suspend/resume):
\begin{lstlisting}
>>> gravando: registro-teste-c-1794329137
>>> API respondendo pelo VIP 51s apos o inicio da retomada
>>> comando vagrant concluido em 80s
OK: registro 'registro-teste-c-1794329137' preservado
\end{lstlisting}

\section{Limitações}
\begin{itemize}
\item \textbf{Tempo do \texttt{vagrant resume}.} O cluster retoma de imediato, mas o
  Vagrant retoma as VMs uma por vez e espera o SSH de cada uma; o comando inteiro leva
  70--80 s. O serviço volta pelo VIP antes disso (51--57 s nas medições), com pouca folga
  em relação aos 60 s do roteiro e dependente do desempenho do host.
\item \textbf{NodePort acessível da sub-rede.} O kube-proxy insere suas regras de DNAT
  antes das do UFW, então a porta 30080 dos nós K3s responde a qualquer host de
  192.168.56.0/24, não só aos proxies. O ponto de entrada documentado é o VIP.
\item \textbf{Banco em um único nó.} Com \texttt{local-path}, o volume do PostgreSQL pertence
  ao nó onde foi criado; se esse nó cair, o banco fica indisponível até ele voltar. Os
  dados ficam no disco da VM e também são copiados em JSON para o host pela VM 4.
\end{itemize}

\section{Como reproduzir}
\begin{lstlisting}
cd tp2
vagrant up
bash scripts/teste-a-balanceamento.sh 10
bash scripts/teste-b-failover.sh
bash scripts/teste-c-ciclo-de-vida.sh          # suspend/resume
bash scripts/teste-c-ciclo-de-vida.sh halt     # halt/up
bash scripts/teste-d-reconciliacao.sh
vagrant provision bastion                      # idempotencia: changed=0
\end{lstlisting}

\end{document}
